\documentclass{article}
\usepackage{graphicx} % Required for inserting images
\usepackage{tikz}
\usepackage{tikz-3dplot}
\usepackage[spanish]{babel}
\usepackage{fontspec}
\usepackage{amsmath, amssymb}
\usepackage{bm}

\title{Formalismo Matemático RA}
\author{Estudiantes RA 2026 2 }
\date{}

\begin{document}

\maketitle

\section{Marcos de referencia y matrices de rotation}

\paragraph{Nombre del estudiante 1}

\paragraph{Nombre del estudiante 1}

\paragraph{Brayan Manrique}

\section*{Inversión Dinámica del Cuadrotor}

\subsection*{Simplificación de la matriz de rotación ($\psi=0$)}

Se asume yaw nulo para reducir la matriz de rotación 3-2-1 completa a una versión que solo depende de $\phi$ y $\theta$, ya que el análisis de inversión dinámica no considera movimiento de guiñada.

\begin{align}
\psi = 0, \quad \cos\psi = 1, \quad \sin\psi = 0
\end{align}

\begin{align}
R^E_B = \begin{bmatrix}
\cos\theta & \sin\phi\sin\theta & \cos\phi\sin\theta \\
0 & \cos\phi & -\sin\phi \\
-\sin\theta & \sin\phi\cos\theta & \cos\phi\cos\theta
\end{bmatrix}
\end{align}

\subsection*{Vector de gravedad proyectado en el marco cuerpo}

Se proyecta el vector de gravedad (definido en el marco inercial NED) hacia el marco cuerpo mediante la matriz de rotación, para obtener cómo ``siente'' el dron la gravedad según su inclinación.

\begin{align}
R^B_E \begin{bmatrix} 0 \\ 0 \\ mg \end{bmatrix} =
\begin{bmatrix} -mg\sin\theta \\ mg\sin\phi\cos\theta \\ mg\cos\phi\cos\theta \end{bmatrix}
\end{align}

\subsection*{Dinámica traslacional en el marco cuerpo}

Se aplica la segunda ley de Newton en cada eje del marco cuerpo: en $X$ e $Y$ solo actúa la componente de gravedad (no hay actuador en esos ejes), mientras que en $Z$ se suma el empuje total de los rotores $F_t$, dividido por la masa.

\begin{align}
\dot u = -g\sin\theta, \qquad
\dot v = g\cos\theta\sin\phi, \qquad
\dot w = g\cos\theta\cos\phi - \frac{F_t}{m}
\end{align}

\subsection*{Proyección a aceleración inercial (NED)}

Se transforma la aceleración traslacional del marco cuerpo al marco inercial mediante $R^E_B$, obteniendo las aceleraciones $V_x, V_y, V_z$ que corresponden a $\ddot p_N, \ddot p_E, \ddot p_D$ --- las variables que maneja el controlador de posición.

\begin{align}
\begin{bmatrix} V_x \\ V_y \\ V_z \end{bmatrix} = R^E_B
\begin{bmatrix} \dot u \\ \dot v \\ \dot w \end{bmatrix}
\end{align}

\subsection*{Expansión algebraica}

Se desarrolla el producto matricial y se simplifica aplicando la identidad pitagórica $\sin^2 x+\cos^2 x=1$, cancelando los términos de gravedad cruzados y dejando expresiones compactas en función únicamente de $\phi$, $\theta$ y $F_t/m$.

\begin{align}
V_x &= -\frac{F_t}{m}\sin\theta\cos\phi \\
V_y &= \frac{F_t}{m}\sin\phi \\
V_z &= g - \frac{F_t}{m}\cos\phi\cos\theta
\end{align}

\subsection*{Inversión dinámica --- despeje del empuje}

Se define $a=F_t/m$ y se eleva al cuadrado cada ecuación, sumándolas: al aplicar la identidad pitagórica dos veces (una para $\theta$, otra para $\phi$), todos los términos angulares se cancelan y queda una expresión directa para la magnitud de aceleración requerida, sin depender de los ángulos.

\begin{align}
a &\equiv \frac{F_t}{m} \\
V_x^2+V_y^2+(g-V_z)^2 &= a^2\left(\sin^2\theta\cos^2\phi+\sin^2\phi+\cos^2\phi\cos^2\theta\right) = a^2 \\
a &= \sqrt{V_x^2+V_y^2+(g-V_z)^2}, \qquad F_t = m\,a
\end{align}

\subsection*{Inversión dinámica --- despeje del roll deseado}

Con $a$ ya conocido, se despeja directamente $\phi_{ref}$ de la ecuación de $V_y$, que depende únicamente de $\phi$.

\begin{align}
\phi_{ref} = \arcsin\left(\frac{V_y}{a}\right)
\end{align}

\subsection*{Inversión dinámica --- despeje del pitch deseado}

Con $a$ y $\phi_{ref}$ ya conocidos, se sustituyen en la ecuación de $V_x$ y se despeja $\theta_{ref}$, cerrando el sistema por sustitución en cascada.

\begin{align}
\theta_{ref} = \arcsin\left(\frac{-V_x}{a\cos\phi_{ref}}\right)
\end{align}

\subsection*{Resumen del bloque de inversión dinámica}

Entradas: aceleraciones deseadas $V_{x,des}, V_{y,des}, V_{z,des}$ provenientes del controlador de posición. Salidas: referencias de actitud $\phi_{ref}, \theta_{ref}$ y empuje total $F_t$, que alimentan los controladores PID de actitud.

\begin{align}
a=\sqrt{V_x^2+V_y^2+(g-V_z)^2}, \quad
F_t = m\,a, \quad
\phi_{ref}=\arcsin\!\left(\frac{V_y}{a}\right), \quad
\theta_{ref}=\arcsin\!\left(\frac{-V_x}{a\cos\phi_{ref}}\right)
\end{align}

\paragraph{Nombre del estudiante 3}

\paragraph{David Santiago Sanabria Romero}
\newcommand{\C}[1]{\cos(#1)}
\newcommand{\Sin}[1]{\sin(#1)}
\section*{Cinemática y Transformación de Coordenadas}

% --- 1. Transformación v2 -> b ---
\subsection*{1. Transformación del marco $v_2$ al marco del cuerpo $b$}
\begin{equation}
\begin{bmatrix}
\hat{x}_{v2} \\
\hat{y}_{v2} \\
\hat{z}_{v2}
\end{bmatrix}
=
\begin{bmatrix}
1 & 0 & 0 \\
0 & \cos\phi & -\sin\phi \\
0 & \sin\phi & \cos\phi
\end{bmatrix}
\begin{bmatrix}
\hat{x}_b \\
\hat{y}_b \\
\hat{z}_b
\end{bmatrix}
\end{equation}

Proyección del eje $\hat{y}_{v2}$:
\begin{equation}
\hat{y}_{v2} = \cos(\phi)\hat{y}_b - \sin(\phi)\hat{z}_b
\end{equation}

\vspace{0.3cm}

% --- 2. Transformación v1 -> b ---
\subsection*{2. Transformación del marco $v_1$ al marco del cuerpo $b$}

Producto de matrices de rotación $R_y(\theta) \cdot R_x(\phi)$:
\begin{equation}
\begin{bmatrix}
\hat{x}_{v1} \\
\hat{y}_{v1} \\
\hat{z}_{v1}
\end{bmatrix}
=
\begin{bmatrix}
\cos\theta & 0 & \sin\theta \\
0 & 1 & 0 \\
-\sin\theta & 0 & \cos\theta
\end{bmatrix}
\begin{bmatrix}
1 & 0 & 0 \\
0 & \cos\phi & -\sin\phi \\
0 & \sin\phi & \cos\phi
\end{bmatrix}
\begin{bmatrix}
\hat{x}_b \\
\hat{y}_b \\
\hat{z}_b
\end{bmatrix}
\end{equation}

Matriz de transformación resultante:
\begin{equation}
\begin{bmatrix}
\hat{x}_{v1} \\
\hat{y}_{v1} \\
\hat{z}_{v1}
\end{bmatrix}
=
\begin{bmatrix}
\cos\theta & \sin\theta\sin\phi & \sin\theta\cos\phi \\
0 & \cos\phi & -\sin\phi \\
-\sin\theta & \cos\theta\sin\phi & \cos\phi\cos\theta
\end{bmatrix}
\begin{bmatrix}
\hat{x}_b \\
\hat{y}_b \\
\hat{z}_b
\end{bmatrix}
\end{equation}

Proyección del eje $\hat{z}_{v1}$:
\begin{equation}
\hat{z}_{v1} = -\sin(\theta)\hat{x}_b + \cos(\theta)\sin(\phi)\hat{y}_b + \cos(\phi)\cos(\theta)\hat{z}_b
\end{equation}

\vspace{0.3cm}

% --- 3. Velocidad Angular ---
\subsection*{3. Vector de Velocidad Angular ($\boldsymbol{\omega}_{I}^b$)}

Suma de las velocidades angulares sobre sus respectivos ejes de rotación intermedios:
\begin{equation}
\begin{aligned}
\bm{\omega}_{I}^b &= \dot{\phi} \hat{x}_b 
+ \dot{\theta} \left( \cos\phi \hat{y}_b - \sin\phi \hat{z}_b \right) \\
&\quad + \dot{\psi} \left( -\sin\theta \hat{x}_b + \cos\theta\sin\phi \hat{y}_b + \cos\phi\cos\theta \hat{z}_b \right)
\end{aligned}
\end{equation}

Agrupando componentes por los ejes del cuerpo ($\hat{x}_b, \hat{y}_b, \hat{z}_b$):
\begin{equation}
\begin{aligned}
\bm{\omega}_{I}^b &= \left( \dot{\phi} - \dot{\psi}\sin\theta \right) \hat{x}_b \\
&\quad + \left( \dot{\theta}\cos\phi + \dot{\psi}\cos\theta\sin\phi \right) \hat{y}_b \\
&\quad + \left( -\dot{\theta}\sin\phi + \dot{\psi}\cos\phi\cos\theta \right) \hat{z}_b
\end{aligned}
\end{equation}

\vspace{0.3cm}

% --- 4. Relación de p, q, r ---
\subsection*{4. Ecuaciones Cinemáticas de Euler}

Componentes de la velocidad angular en los ejes fijos al cuerpo ($p, q, r$):
\begin{align}
p &= \dot{\phi} - \dot{\psi}\sin\theta \\
q &= \dot{\theta}\cos\phi + \dot{\psi}\cos\theta\sin\phi \\
r &= -\dot{\theta}\sin\phi + \dot{\psi}\cos\phi\cos\theta
\end{align}

Expresado en forma matricial:
\begin{equation}
\begin{bmatrix}
p \\
q \\
r
\end{bmatrix}
=
\begin{bmatrix}
1 & 0 & -\sin\theta \\
0 & \cos\phi & \cos\theta\sin\phi \\
0 & -\sin\phi & \cos\phi\cos\theta
\end{bmatrix}
\begin{bmatrix}
\dot{\phi} \\
\dot{\theta} \\
\dot{\psi}
\end{bmatrix}
\end{equation}

Por lo tanto, la tasa de variación de los ángulos de Euler es función de las velocidades angulares del cuerpo:
\begin{equation}
(\dot{\phi}, \dot{\theta}, \dot{\psi}) = f (p, q, r)
\end{equation}

\paragraph{Profe}

\subsection{Dinámica traslacional en el marco del cuerpo}

\subsubsection{Segunda ley de Newton}

\begin{equation}
 \left(m \frac{d\vec{V}}{dt} \right)_w = \sum \vec{F}
\end{equation}

Partiendo de la segunda ley de Newton en el marco rotante:

\begin{equation}
m \left( \frac{d\vec{V}}{dt} \right)_b + m \, \vec{\boldsymbol{\omega}} \times \vec{V}_b = \vec{F}
\end{equation}

Expresando en forma matricial:

\begin{equation}
m
\begin{bmatrix}
\dot{u} \\
\dot{v} \\
\dot{w}
\end{bmatrix}
+
m
\begin{bmatrix}
\hat{i} & \hat{j} & \hat{k} \\
p & q & r \\
u & v & w
\end{bmatrix}
=
\begin{bmatrix}
F_x \\
F_y \\
F_z
\end{bmatrix}
\end{equation}

\paragraph{Javier Alejandro Tellez Paez}

\section*{Dinámica}

Ecuación general de transporte (derivada de un vector en marco inercial vs.\ body):
\begin{equation}
\left(\frac{d\vec{A}}{dt}\right)_{I}=\left(\frac{d\vec{A}}{dt}\right)_{b}+\vec{\omega}\times\vec{A}
\end{equation}

Momento lineal $\vec{P}=m\vec{V}$ y momento angular $\vec{H}=\mathbf{I}\vec{\omega}$. Con $m$ constante:
\begin{align}
\frac{d\vec{P}}{dt}&=\frac{d(m\vec{V})}{dt}=m\,\frac{d\vec{V}}{dt}=\sum \vec{F}\\[4pt]
\frac{d(\mathbf{I}\vec{\omega})}{dt}&=\mathbf{I}\,\frac{d\vec{\omega}}{dt}=\sum \vec{\tau}
\end{align}

Aplicado al problema, con $(\cdot)_{I}$ marco inercial y $(\cdot)_{b}$ marco body:
\begin{align}
m\left(\frac{d\vec{V}}{dt}\right)_{I}&=\sum \vec{F}\\[4pt]
m\left[\left(\frac{d\vec{V}}{dt}\right)_{b}+\vec{\omega}\times\vec{V}_{b}\right]&=\sum \vec{F}
\end{align}

En componentes, con $\vec{\omega}=(p,q,r)$ y $\vec{V}_{b}=(u,v,w)$:
\begin{equation}
m\begin{bmatrix}\dot{u}\\[2pt]\dot{v}\\[2pt]\dot{w}\end{bmatrix}
+m\begin{vmatrix}\hat{\imath}&\hat{\jmath}&\hat{k}\\ p&q&r\\ u&v&w\end{vmatrix}
=\begin{bmatrix}F_{x}\\[2pt]F_{y}\\[2pt]F_{z}\end{bmatrix}
\end{equation}

Desarrollando el producto cruz:
\begin{equation}
m\begin{bmatrix}\dot{u}\\[2pt]\dot{v}\\[2pt]\dot{w}\end{bmatrix}
+m\begin{bmatrix}qw-rv\\[2pt]ru-pw\\[2pt]pv-qu\end{bmatrix}
=\begin{bmatrix}F_{x}\\[2pt]F_{y}\\[2pt]F_{z}\end{bmatrix}
\end{equation}

Despejando las aceleraciones:
\begin{equation}
\begin{bmatrix}\dot{u}\\[2pt]\dot{v}\\[2pt]\dot{w}\end{bmatrix}
=\frac{1}{m}\begin{bmatrix}F_{x}\\[2pt]F_{y}\\[2pt]F_{z}\end{bmatrix}
-\begin{bmatrix}qw-rv\\[2pt]ru-pw\\[2pt]pv-qu\end{bmatrix}
\end{equation}

Forma final para simulación:
\begin{equation}
\boxed{
\begin{bmatrix}\dot{u}\\[2pt]\dot{v}\\[2pt]\dot{w}\end{bmatrix}
=\frac{1}{m}\begin{bmatrix}F_{x}\\[2pt]F_{y}\\[2pt]F_{z}\end{bmatrix}
+\begin{bmatrix}rv-qw\\[2pt]pw-ru\\[2pt]qu-pv\end{bmatrix}}
\end{equation}




\end{document}
